\section{Мета}
\label{sec:objective}
Метою даної лабораторної роботи є дослідження принципів побудови сучасних електронних платіжних систем, 
криптографічних механізмів, що забезпечують захист платіжної операції на кожному з її етапів та основних загроз, 
що виникають під час здійснення електронних платежів.

Для досягнення цієї мети було поставлено такі завдання:
\begin{enumerate}[noitemsep]
    \item Розглянути основні механізми безпеки електронних карткових платежів (TLS, токенізація, \\ CVV/CVC, 3-D 
        Secure, strong client authentication, цифрові підписи та HMAC, PCI DSS) з використанням чинних стандартів та технічної документації;
    \item Побудовати та проаналізувати поетапно спрощену схему онлайн-платежу карткою;
    \item Проаналізувати п'ять типових ситуацій порушення безпеки разом із відповідними заходами протидії для 
        кожної з них;
    \item Дослідити сучасні спосіб оплати банківський переказ, ініційований через Open Banking;
    \item Розглянути історичну схему мікроплатежів MicroMint та провести її порівняння з сучасними механізмами 
        безпеки платежів.
\end{enumerate}


\section{Механізми безпеки сучасних електронних платежів}
\label{sec:questions}



\paragraph{Питання 1. Учасники онлайн-платежу за допомогою картки та їхні ролі.} У платіжній операції за допомогою 
картки завжди беруть участь однакові сторони, навіть якщо клієнт бачить лише сторінку оформлення замовлення продавця~\cite{stripe_card_process}:
\begin{itemize}[noitemsep]
    \item \textbf{Клієнт (cardholder)} -- володіє карткою та рахунком, з якого знімаються кошти. Клієнт вводить 
        реквізити картки (або обирає збережену картку / електронний гаманець) і, якщо це потрібно, підтверджує 
        платіж.
    \item \textbf{Продавець (merchant)} -- інтернет-магазин, що продає товари або послуги. Він створює замовлення, 
        надсилає запит на оплату своєму платіжному провайдеру та, зрештою, отримує гроші.
    \item \textbf{Платіжний провайдер (PSP, Payment Service Provider) / acquirer} -- еквайєр (банк-еквайєр) 
        обробляє карткові платежі від імені продавця та направляє їх до карткової мережі; PSP або платіжний процесор 
        (наприклад, Stripe) є технічним посередником, який часто поєднує обидві функції~\cite{stripe_card_process}.
    \item \textbf{Платіжна система (card network)} -- Visa, Mastercard тощо. Вона направляє запит на авторизацію 
        до відповідного емітента за номером картки, а згодом здійснює переказ коштів між банками під час клірингу 
        та розрахунків~\cite{stripe_card_process}.
    \item \textbf{Емітент (issuing bank)} --банк, який випустив картку. Він здійснює аутентифікацію власника картки 
            (наприклад, через 3-D Secure), перевіряє наявність коштів, ліміти та ознаки шахрайства, а також 
            \emph{схвалює або відхиляє} транзакцію~\cite{stripe_card_process}.
\end{itemize}
\textit{Приклад:} Клієнт купує книгу за 600 грн. Магазин (торговець) користується послугами PSP; еквайєр PSP 
надсилає запит через Visa (мережу) до банку клієнта (емітента), який резервує 600 грн і повертає код схвалення.

\paragraph{Питання 2. Ідентифікація користувача vs авторизації платежу.}
\textbf{Ідентифікація (Authentication)} відповідає на питання \enquote{\emph{Хто ви?}} -- система перевіряє, чи 
дійсно особа є законним власником картки або рахунку, наприклад, за допомогою пароля, одноразового коду або 
відбитка пальця~\cite{nist_80063b4}. \textbf{Авторизація (Authorisation)} платежу відповідає на запитання 
\enquote{\emph{Чи дозволена ця конкретна операція?}} -- емітент перевіряє, чи картка дійсна та не заблокована, 
чи є достатньо коштів або кредиту, і чи збігаються дані; потім він надсилає код схвалення або відмови та 
резервує суму платежу~\cite{stripe_card_process}.

Аутентифікація стосується \emph{особи}, авторизація -- \emph{транзакції}. Користувачеві, який пройшов повну 
аутентифікацію, все одно може бути відмовлено (через недостатність коштів), а платіж може бути авторизовано 
без етапу посиленої аутентифікації, якщо застосовується виняток (наприклад, платіж на невелику суму, більше про 
це в питані \No 7).

\paragraph{Питання 3. Роль протоколу TLS (Transport Layer Security) в електронному платежі.}
TLS створює захищений канал між двома кінцевими точками -- браузером клієнта та інтернет-магазином, інтернет-магазином 
та PSP, PSP та еквайром. Його поточна версія, TLS 1.3, була опублікована в липні 2026 року як RFC~9846, що 
замінює RFC~8446~\cite{rfc9846}. TLS забезпечує три властивості~\cite{rfc9846}:

\begin{itemize}[noitemsep]
    \item \textbf{Автентифікація сервера} -- браузер перевіряє сертифікат сервера, тому зловмисник не може 
        видати себе за інтернет-магазин (захист від атак типу MiTM та фішингу на мережевому рівні, за умови, 
        що користувач перебуває на правильному домені);
    \item \textbf{Конфіденційність} -- дані картки, файли cookie та деталі замовлення шифруються, тому пасивне 
        перехоплення інформації є марним;
    \item \textbf{Цілісність} -- кожен запис захищений за допомогою AEAD (Authenticated Encryption with Associated 
        Data), тому будь-які зміни під час передачі виявляються.
\end{itemize}
TLS \emph{не приховує} обсяг даних або сам факт зв'язку~\cite{rfc9846} і захищає дані \emph{лише під час передачі}: 
щойно дані досягають сервера, TLS більше не допомагає (див. розділ~\ref{sec:tls-not-enough}).

\paragraph{Питання 4. Токенізація платіжних даних.}
Токенізація замінює конфіденційний номер картки -- PAN (Primary Account Number) -- на substitute value, тобто 
платіжний токен. Відповідність "токен $\leftrightarrow$ PAN"{} зберігається лише в захищеному сховищі токенів 
постачальника токенів (TSP) і може бути відновлена лише уповноваженими суб'єктами~\cite{emv_tokenisation}. 
Платіжний токен EMV додатково прив'язаний за допомогою \emph{Token Domain Restriction Controls} до конкретного продавця, 
пристрою або каналу, тому токен, викрадений в одному магазині, не може бути використаний деінде~\cite{emv_tokenisation_qrg}. 
Переваги зберігання токена замість PAN:

\begin{itemize}[noitemsep]
    \item У разі витоку даних отримані дані є марними поза межами їхньої домени~\cite{emv_tokenisation_qrg};
    \item Втрачений пристрій або скомпрометований токен можна \emph{від'єднати} без перевипуску фізичної 
        картки~\cite{emv_tokenisation};
    \item Торговець не зберігає номери PAN, що суттєво зменшує сферу застосування стандарту PCI DSS scope (Q9).
\end{itemize}
В Україні термін \emph{card token} було додано до Положення НБУ про аутентифікацію у 2026 році: токен, створений 
після сильної аутентифікації емітентом, може сам по собі слугувати одним із елементів SCA (\enquote{possession})~\cite{nbu4_2026}. 
\textit{Приклад:} Apple Pay зберігає токен \texttt{4895\,xxxx\,xxxx\,1234} на телефоні; навіть якщо база даних 
магазину витокує його, токен працює лише для Apple Pay на цьому конкретному пристрої.

\paragraph{Питання 5. Призначення CVV/CVC та причини обмеження його зберігання}
CVV2/CVC2 (Card Verification Value / Card Validation Code) -- це трьо/чотиризначний код захисту картки. Його 
запитують під час транзакцій без фізичної присутності картки (card-not-present, CNP), щоб підтвердити, що клієнт 
\emph{фізично володіє} карткою, а не лише знає її номер~\cite{pci_faq_cvv_request}. Оскільки цей код не закодований 
на магнітній смузі чи в чіпі, викрадена база даних номерів карток без них стає набагато менш корисною. Саме тому 
стандарт PCI DSS класифікує цей код як \emph{Sensitive Authentication Data (SAD)} і забороняє його зберігання 
після авторизації -- навіть у зашифрованому вигляді та навіть за згодою клієнта~\cite{pci_faq_1210, pci_faq_1280}. 
Якщо торговець зберігав би CVV-коди, одного витоку даних було б достатньо, щоб надати зловмисникам усе необхідне 
для здійснення онлайн-шахрайства.

\paragraph{Питання 6. Мета 3-D Secure та додаткової автентифікації користувача.}
3-D Secure (3DS, \enquote{three-domain secure}: домен merchant/acquirer, домен взаємодії та домен емітента) -- це 
протокол EMVCo, який дозволяє торговцю та емітенту обмінюватися даними з метою автентифікації власника картки під 
час онлайн-покупки та запобігання шахрайству CNP~\cite{emv_3ds}. Сервер 3DS продавця надсилає дані про транзакцію, 
платіж та пристрій через сервер каталогів карткової мережі на сервер контролю доступу (Access Control Server, ACS) 
емітента~\cite{emv_3ds_whitepaper}. На основі цих даних емітент приймає рішення:
\begin{itemize}[noitemsep]
    \item \textbf{Безперешкодний потік (frictionless flow)} -- низький ризик, клієнт просто натискає \enquote{Buy};
    \item \textbf{Процес із перевіркою (challenge flow)} -- вищий ризик, клієнт повинен підтвердити операцію за 
        допомогою одноразового пароля, біометричних даних або іншим способом~\cite{emv_3ds}.
\end{itemize}
Результатом є криптографічне значення автентифікації, яке додається до запиту на авторизацію та підтверджує,
що емітент автентифікував власника картки. Таким чином, 3DS є практичним способом реалізації двофакторної
сильної автентифікації клієнта (strong customer authentication, SCA) для карткових платежів~\cite{emv_3ds}.

\paragraph{Вимоги до сильної автентифікації в Україні (Постанова НБУ \No~58)}
\emph{Положення про автентифікацію та застосування сильної автентифікації на платіжному ринку} (Постанова НБУ 
\No~58 від 3 травня 2023 року, із змінами, внесеними Постановою \No~4 від 13 січня 2026 року) встановлює такі 
основні правила~\cite{nbu58, nbu4_2026}:
\begin{itemize}[noitemsep]
    \item \textbf{Означення.} Посилена автентифікація використовує щонайменше два елементи з \emph{різних} категорій: 
        знання (те, що знає лише користувач), володіння (предмет, який має лише користувач) та властивість 
        (біометричні дані)~\cite[п.~3(15)]{nbu58}. Компрометація одного елемента не повинна становити загрози 
        для інших~\cite[п.~32]{nbu58}.
    \item \textbf{Коли це є обов'язковим.} Для віддаленого доступу до облікового запису, для ініціювання віддаленого 
        платежу та для будь-якої дії, щодо якої є підозра у шахрайстві~\cite[п.~7, 15]{nbu58}.
    \item \textbf{Код автентифікації.} Результатом SCA є код автентифікації, який повинен бути захищеним від підробки, 
        не повинен розкривати жодних елементів і не повинен бути виведеним із попередніх кодів; він реалізується 
        за допомогою одноразових паролів, (кваліфікованих) електронних підписів або інших криптографічних засобів~\cite[п.~14, 16]{nbu58}, \cite{nbu4_2026}.
    \item \textbf{Динамічне зв'язування (linking).} Платник бачить суму та отримувача; код прив'язаний саме до цієї 
        суми та отримувача, і будь-яка зміна цих даних робить код недійсним~\cite[п.~22]{nbu58}.
    \item \textbf{Безпека сеансу.} Не більше п'яти послідовних невдалих спроб, після чого аутентифікація блокується; 
        не більше десяти хвилин бездіяльності після входу в систему~\cite[п.~17]{nbu58}.
    \item \textbf{Винятки (Exemptions).} SCA можна пропустити, наприклад, для дистанційних платежів на суму до 
    2\,000~UAH (з сукупними лімітами 10\,000~UAH або п'ятьма послідовними операціями), для довірених одержувачів, 
    для наступних платежів у рамках серії регулярних платежів, для переказів між власними рахунками та для операцій 
    з низьким рівнем ризику згідно з аналізу ризиків транзакцій (transaction risk analysis, TRA) з орієнтовними 
    показниками шахрайства для порогових значень 2\,000, 7\,000 і 20\,000~UAH~\cite[п.~37--47, Додаток]{nbu58}. 
    З 2026 року також звільняються від SCA податки та комунальні платежі на суму до 30\,000~UAH~\cite{nbu4_2026}.
\end{itemize}
\textit{Приклад:} Для переказу 5\,000~UAH новому одержувачу в банківському додатку потрібно, наприклад, ввести 
PIN-код (knowledge) на зареєстрованому смартфоні (knowledge), а на екрані підтвердження відображається точна сума 
та одержувач.

\paragraph{Питання 8. Цифрові підписи та HMAC між компонентами платіжної системи.}
Протокол TLS захищає повідомлення лише між двома сусідніми вузлами. Платіжні повідомлення часто проходять через 
кілька систем і реєструються, ставляться в чергу або пересилаються, тому їхня цілісність та походження мають 
бути перевірені на всьому шляху (\emph{end-to-end}):
\begin{itemize}[noitemsep]
    \item \textbf{HMAC (Hash-based Message Authentication Code)} обчислює тег на основі повідомлення та спільного 
        секретного ключа за допомогою геш-функції\cite{rfc2104}. Це доводить, що повідомлення було створено особою, 
        яка знає ключ, і не зазнало змін. \textit{Приклад:} Stripe підписує кожен вебхук (наприклад, \enquote{платіж 
        успішно здійснено}) за допомогою HMAC-SHA256 на основі \texttt{timestamp.body}; магазин перераховує тег і відхиляє 
        підроблені або застарілі сповіщення~\cite{stripe_webhook_signature}.
    \item \textbf{Digital signature} використовує приватний ключ для підписання та відкритий ключ для перевірки. 
        Окрім цілісності, він забезпечує незаперечність (non-repudiation): підписувач не може згодом заперечити факт 
        надсилання  повідомлення, оскільки ніхто інший не володіє приватним ключем. НБУ вимагає використання електронних 
        підписів, унікальних ідентифікаторів сеансу та міток часу за UTC у електронній взаємодії між учасниками платіжного 
        ринку~\cite[п.~74--75]{nbu58}, а в рамках Open Banking банк може вимагати від сервісу ініціювання платежів 
        завіряти свої повідомлення кваліфікованою електронною печаткою~\cite[п.~30]{nbu80}.
\end{itemize}

\paragraph{Мета стандарту PCI DSS та що він захищає.}
PCI DSS (Payment Card Industry Data Security Standard), який підтримується Радою з питань стандартів безпеки PCI, 
є набором базових технічних та операційних вимог, призначених для захисту даних платіжних рахунків~\cite{pci_dss}. 
Поточна версія -- v4.0.1. Він захищає дві групи \emph{account data}:
\begin{itemize}[noitemsep]
    \item \textbf{Дані власника картки (Cardholder Data, CHD)} -- номер платіжної картки (PAN), ім'я власника картки, 
        термін дії, код послуги; їх можна зберігати лише за умови захисту (наприклад, PAN має бути зроблений нечитабельним 
        шляхом шифрування, обрізання або токенізації);
    \item \textbf{Чутливі дані автентифікації (Sensitive Authentication Data, SAD)} -- повні дані магнітної смуги, 
        коди перевірки картки та PIN-коди/блоки PIN-кодів; їх ніколи не можна зберігати після авторизації~\cite{pci_faq_1210}.
\end{itemize}
Стандарт застосовується до всіх суб'єктів, які зберігають, обробляють або передають CHD/SAD або можуть впливати на 
безпеку середовища даних власника картки (Cardholder Data Environment, CDE): торговців, процесорів, еквайрів, 
емітентів та постачальників послуг~\cite{pci_dss}. Він охоплює всю інфраструктуру: засоби контролю мережевої 
безпеки, безпечну конфігурацію, шифрування даних у стані спокою та під час передачі, управління вразливостями, 
контроль доступу та багатофакторна автентифікація (Multi-Factor Authentication, MFA), ведення журналів та 
моніторинг, регулярне тестування та політики безпеки.

\section{Карткові онлайн-розрахунки}
\label{sec:flow}
Розглянемо клієнта, який купує товари в інтернет-магазині та розраховується банківською карткою. Магазин
використовує сучасного PSP: поля для введення даних картки на сторінці оформлення замовлення
обслуговуються безпосередньо PSP (hosted fields), тому власний сервер магазину ніколи не бачить повного
номера картки. На рисунку~\ref{fig:card-flow} показано спрощений алгоритм дій.

\begin{figure}[H]
    \centering
    \begin{tikzpicture}[node distance = 1.15cm]
        \tikzset{actor/.append style = {minimum width = 2.2cm, text width = 2.0cm}}
        \node[actor] (cust) {Customer\\(browser)};
        \node[actor, right = of cust] (merch) {Merchant\\(online shop)};
        \node[actor, right = of merch] (psp) {PSP /\\Acquirer};
        \node[actor, right = of psp] (net) {Card network\\(DS)};
        \node[actor, right = of net] (iss) {Issuer bank\\(ACS)};

        % requests (top)
        \draw[flow] ([yshift = 3mm]cust.east) -- node[above, lbl] {(2)\\order} ([yshift = 3mm]merch.west);
        \draw[flow] ([yshift = 3mm]merch.east) -- node[above, lbl] {(2) token,\\amount} ([yshift = 3mm]psp.west);
        \draw[flow] ([yshift = 3mm]psp.east) -- node[above, lbl] {(4) auth.\\request} ([yshift = 3mm]net.west);
        \draw[flow] ([yshift = 3mm]net.east) -- node[above, lbl] {(4) routed\\by BIN} ([yshift = 3mm]iss.west);

        % responses (bottom)
        \draw[back] ([yshift = -3mm]iss.west) -- node[below, lbl] {(5) approve\\/ decline} ([yshift = -3mm]net.east);
        \draw[back] ([yshift = -3mm]net.west) -- node[below, lbl] {(5)\\response} ([yshift = -3mm]psp.east);
        \draw[back] ([yshift = -3mm]psp.west) -- node[below, lbl] {(6) signed\\webhook} ([yshift = -3mm]merch.east);
        \draw[back] ([yshift = -3mm]merch.west) -- node[below, lbl] {(6)\\receipt} ([yshift = -3mm]cust.east);

        % card data entry directly to PSP
        \draw[flow, dashed] (cust.north) to[out = 35, in = 145]
        node[above, lbl, pos = 0.5, font = \small] {(1) card data via PSP hosted fields $\rightarrow$ token} (psp.north);

        % 3DS challenge
        \draw[flow, warnOrange, dashed] (cust.south) to[out = -25, in = -155]
        node[below, lbl, pos = 0.5, text = warnOrange!70!black]
        {(3) 3-D Secure challenge between customer and issuer ACS: OTP / push / biometrics (SCA)} (iss.south);
    \end{tikzpicture}
    \caption{ Спрощена схема карткового онлайн-платежу з токенізацією та 3-D Secure.}
    \label{fig:card-flow}
\end{figure}

\noindent Основні етапи: (1)~клієнт вводить дані картки у поля, розміщені на сервері PSP, та отримує токен; (2)~замовлення 
та токен надсилаються торговцю, який звертається до PSP із запитом на списання коштів; (3)~за допомогою 3-D Secure 
здійснюється автентифікація власника картки в емітента; (4)~запит на авторизацію надходить через еквайра та 
платіжну мережу до емітента; (5)~емітент схвалює або відхиляє запит; (6)~результат повертається торговцю у вигляді 
відповіді API та підписаного вебхука. Пізніше в процесі клірингу та розрахунків фактичні кошти перераховуються 
від емітента до еквайра та продавця~\cite{stripe_card_process}. У таблиці~\ref{tab:stages} нижче наведено підсумкові 
дані щодо інформації та захисту на кожному етапі.

\begin{longtblr}[
        caption = {Дані, доступ та захист на основних етапах онлайн-платежу карткою.},
        label = {tab:stages},
    ]{
        colspec = {Q[l,m,0.12\linewidth] X[l,m] X[l,m] X[l,m] X[l,m]},
        row{1} = {font = \bfseries, bg = headBlue!15},
        rowhead = 1,
        hlines, vlines,
    }
    Етап & Передані дані & Хто має доступ & Що необхідно захищати & Механізми \\
    1. Введення даних картки &
    PAN, термін дії, CVV2, ім'я власника картки &
    Клієнт, PSP (сервер торговця — ні) &
    Усі поля: конфіденційність; CVV2 не повинен зберігатися &
    TLS~1.3, поля, розміщені на сервері PSP / iframe, негайна токенізація, PCI DSS~\cite{pci_dss} \\
    2. Оформлення замовлення &
    Ідентифікатор замовлення, сума, валюта, токен, ключ ідемпотентності &
    Клієнт, торговець, PSP &
    Сума та одержувач платежу: цілісність; токен: конфіденційність &
    TLS, обчислення суми на стороні сервера, секретний ключ API, ключ ідемпотентності~\cite{stripe_idempotency} \\
    3. 3-D Secure &
    Дані про транзакції та пристрої, одноразовий пароль (OTP) / підтвердження через push-повідомлення, значення автентифікації &
    Сервер 3DS торговця, сервер каталогів, ACS емітента, клієнт &
    Одноразовий пароль (OTP) та біометричні дані: конфіденційність; сума та одержувач платежу, що відображаються користувачеві: цілісність &
    TLS на всіх етапах, SCA з динамічним зв'язуванням, криптографічна автентифікація значення~\cite{emv_3ds, nbu58} \\
    4. Запит на авторизацію &
    PAN або мережевий токен, сума, ідентифікатор торговця, значення автентифікації, криптограма &
    Еквайєр, платіжна мережа, емітент &
    PAN: конфіденційність; сума, торговець: цілісність &
    Зашифровані приватні мережі, MAC-код у повідомленнях, HSM (апаратний модуль безпеки) для ключів, зіставлення токенів із PAN у сховищі~\cite{emv_tokenisation} \\
    5. Рішення &
    Код схвалення / відмови, код авторизації &
    Емітент, мережа, еквайєр, PSP &
    Цілісність відповіді (підробка статусу \enquote{затверджено} має бути неможливою) &
    Автентифіковані канали, MAC, зіставлення за ідентифікатором транзакції \\
    6. Повідомлення &
    Статус платежу, ідентифікатор замовлення, сума &
    ПСП, торговець &
    Цілісність та автентичність статусу &
    Підпис HMAC-SHA256 підпис із позначкою часу на веб-хуках~\cite{stripe_webhook_signature} \\
\end{longtblr}

\textbf{Де і як відбувається автентифікація користувача.} Автентифікацію клієнта здійснює \emph{issuer}, а не 
магазин: під час процедури 3-D Secure система ACS емітента або автентифікує клієнта у фоновому режимі на основі 
даних про ризики (без додаткових дій, frictionless), або вимагає підтвердження -- зазвичай у вигляді push-повідомлення 
в банківському додатку з введенням PIN-коду або відбитка пальця, що є двофакторною SCA~\cite{emv_3ds, nbu58}. 
Інтернет-магазин отримує лише результат, але ніколи не дізнається банківські реквізити користувача.

\textbf{Як операція схвалюється або відхиляється.} Остаточне рішення приймає емітент під час авторизації: він 
перевіряє стан картки (чи не загублена або заблокована), наявність коштів або кредитного ліміту, результат 3DS 
та власний рейтинг ризику шахрайства. Він повертає код відповіді; у разі схвалення сума резервується на рахунку 
та видається код авторизації~\cite{stripe_card_process}. Продавець відправляє товар лише після отримання підтвердженого 
статусу \enquote{успішно}.

\subsection*{Чому одного TLS недостатньо}
\label{sec:tls-not-enough}
TLS є необхідним, але охоплює лише один клас загроз -- атаки на дані \textit{під час передачі (in transit)} між 
двома сусідніми вузлами. Він не вирішує таких проблем:
\begin{itemize}[noitemsep]
    \item \textbf{Дані у стані спокою (Data at rest)} Після розшифрування на сервері дані перебувають у вигляді 
        відкритого тексту; злам бази даних, дії інсайдера або шкідливе програмне забезпечення на сервері не 
        підпадають під дію TLS. Цю проблему вирішують токенізація та стандарт PCI DSS.
    \item \textbf{Ендпоїни (Endpoints)} Шкідливе програмне забезпечення (Malware) в браузері клієнта (наприклад, 
        скімер, вбудований у сторінку оформлення замовлення) зчитує дані \textit{ще до} їх шифрування.
    \item \textbf{Ідентичність платника} TLS перевіряє автентичність \emph{сервера}, а не власника картки. Злодій 
        із викраденими даними картки може встановити цілком дійсний сеанс TLS; тут необхідні SCA/3-D Secure.
    \item \textbf{Фішинг.} Фальшивий інтернет-магазин може мати дійсний сертифікат для власного домену, схожого 
        на справжній.
    \item \textbf{Бізнес-логіка} Повторне відтворення правильного запиту, подвійне списання коштів або зміна 
        суми, спричинені легітимним, але скомпрометованим компонентом, не виявляються за допомогою TLS; необхідні 
        ідемпотентні ключі, нонси, HMAC/підписи та динамічне зв'язування.
    \item \textbf{Ланцюги з декількома проміжними вузлами (Multi-hop chains)} TLS працює "від вузла до вузла"{}(hop-by-hop): 
        платіжний сервіс-провайдер (PSP), мережа та емітент по черзі завершують сеанс, тому для забезпечення цілісності 
        "від кінця до кінця"{} та неможливості заперечення авторства потрібні підписи на рівні повідомлень.
\end{itemize}

\section{Аналіз ситуацій порушення безпеки}
\label{sec:situations}

\subsection{Ситуація 1: Перехоплення мережевого трафіку}
\label{sec:s1}

Зловмисник може спостерігати за трафіком між клієнтом та магазином; з'єднання правильно використовує
HTTPS/TLS (Малюнок~\ref{fig:s1}).

\begin{figure}[H]
    \centering
    \begin{tikzpicture}
        \node[actor] (c) {Customer};
        \node[actor, right = 9cm of c] (m) {Online shop};
        \draw[very thick, okGreen, double, double distance = 4pt] (c) --
        node[above = 4pt, lbl, text = okGreen!60!black] {\textbf{TLS 1.3 tunnel}: card data, cookies,
        order, amount\\(encrypted + integrity-protected)} (m);
        \node[evil, below = 1.4cm of $(c)!0.5!(m)$] (e) {Eavesdropper};
        \draw[atk] ($(c)!0.5!(m)$) -- (e);
        \node[note, right = 0.4cm of e, text width = 4.6cm] {
            \textcolor{badRed}{\textbf{Visible:}} IP addresses, shop domain (SNI, DNS),
            time and size of packets\\
            \textcolor{okGreen}{\textbf{Hidden:}} URL path, card data, CVV, password, cookies, amount};
    \end{tikzpicture}
    \caption{Situation~1: passive observation of a correctly configured TLS connection.}
    \label{fig:s1}
\end{figure}

\textbf{Що може дізнатися зловмисник.} Лише метадані: хто з ким спілкується (IP-адреси), ім'я домену (зазвичай 
видно з Server Name Indication та DNS-запитів), коли та скільки даних надсилається. TLS явно не приховує довжину 
повідомлень~\cite{rfc9846}, тому аналіз трафіку може виявити, наприклад, чи була завантажена сторінка оплати.

\textbf{Від чого захищає TLS.} Зчитування вмісту (дані картки, CVV, логін, сесійні файли cookie), модифікацію 
або injecting пакетів (будь-яка зміна порушує тег AEAD), видавання себе за інтернет-магазин (перевірка сертифіката 
завершиться невдачею) та повторне відтворення записаних записів (послідовні номери та свіжі сесійні ключі). 
Оскільки TLS~1.3 використовує ефемерний алгоритм Діффі–Хелмана, записаний трафік залишається захищеним навіть у 
випадку подальшого витоку довгострокового ключа сервера (forward secrecy)~\cite{rfc9846}.

\textbf{Наслідки та залишкові ризики.} При правильному використанні TLS наслідки обмежуються витоком метаданих, 
що стосуються конфіденційності. Ризик зростає лише у разі неправильної конфігурації TLS (застарілі версії, вимкнена 
перевірка сертифікатів) або якщо користувач приймає попередження про сертифікат. Заходи протидії: виключно 
TLS~1.3, HSTS (HTTP Strict Transport Security) для запобігання переходу на HTTP, правильна перевірка сертифікатів 
у мобільних додатках.

\subsection{Ситуація 2: Компрометація бази даних продавця}
\label{sec:s2}

Зловмисник отримав копію бази даних інтернет-магазину. На рис.~\ref{fig:s2} порівняно дві стратегії зберігання даних.

\begin{figure}[H]
    \centering
    \begin{tikzpicture}
        \node[db, draw = badRed, fill = badRed!8] (d1) at (0, 0)
        {\textbf{Variant A}\\PAN 4149\ldots0187\\exp. 09/29\\(CVV2 --- forbidden)};
        \node[db, draw = okGreen, fill = okGreen!8] (d2) at (9, 0)
        {\textbf{Variant B}\\\texttt{pm\_9Fq2\ldots}\\\texttt{pm\_K7aZ\ldots}\\last4 = 0187};
        \node[evil] (e1) at (0, -2.7) {Attacker};
        \node[evil] (e2) at (9, -2.7) {Attacker};
        \draw[atk] (d1) -- (e1);
        \draw[atk] (d2) -- (e2);
        \node[note, right = 0.3cm of e1, text width = 4.1cm, draw = badRed]
        {Card numbers usable on other sites, card fraud, mass card reissue, PCI DSS fines};
        \node[note, right = 0.3cm of e2, text width = 3.8cm, draw = okGreen]
        {Tokens are valid only for this merchant's PSP account and require the merchant's secret API key};
    \end{tikzpicture}
    \caption{Situation~2: database leak with stored card data (A) vs.\ stored tokens (B).}
    \label{fig:s2}
\end{figure}

\textbf{Варіант А -- магазин зберігав дані карток.} Зловмисник отримує номери карток (PAN), дати закінчення 
терміну їх дії, а також імена всіх клієнтів. Ці дані можна продати або використати в інших інтернет-магазинах 
(особливо тих, що мають слабкі перевірки 3DS). Якщо магазин також незаконно зберігав коди CVV, витік стає повним 
"набором"{} для шахрайства CNP~\cite{pci_faq_1280}. Наслідки: прямі збитки для клієнтів, масове переоформлення 
карток банками, штрафи та можлива втрата права приймати картки продавцем, втрата репутації.

\textbf{Варіант Б — магазин зберігав лише платіжні токени.} Токени є сурогатними значеннями. Вони обмежені 
доменом: токен може використовуватися лише цим торговцем через його PSP, а PSP додатково вимагає секретного API-ключа 
торговця~\cite{emv_tokenisation_qrg}. За межами цієї доменної зони токен не має вартості, і постачальник токенів 
(TSP) може анулювати або відв'язати токени без перевипуску карток~\cite{emv_tokenisation}. Витік все одно розкриває 
особисті дані (імена, адреси, останні чотири цифри), але не дані карток, придатні для використання. 

\textbf{Роль токенізації.} Токенізація переміщує цінні дані (PAN) із системи продавця у спеціалізоване, захищене 
сховище токенів. Таким чином вона зменшує радіус ураження (blast radius) порушення та обсяг вимог PCI DSS для продавця. Це 
необхідно поєднувати із захистом ключів API: якщо зловмисник викраде базу даних і секретний ключ, він зможе ініціювати 
списання коштів за допомогою токенів -- тому ключі мають зберігатися в окремому сховищі секретних даних, а не в 
тій самій базі даних.

\subsection{Ситуація 3: Повторне виконання операції}
\label{sec:s3}

Зловмисник отримав правильний запит на оплату та повторно надсилає його до платіжної системи (риc.~\ref{fig:s3}).

\begin{figure}[H]
    \centering
    \begin{tikzpicture}
        \node[actor] (m) {Merchant};
        \node[actor, right = 6.8cm of m] (p) {Payment\\system};
        \node[evil, below = 1.8cm of $(m)!0.5!(p)$] (e) {Attacker};
        \draw[flow] ([yshift = 2mm]m.east) -- node[above, lbl]
        {(1) \texttt{POST /charge}: amount = 600, key = \texttt{a81f\ldots},\\ts, nonce, HMAC}
        ([yshift = 2mm]p.west);
        \draw[atk] ($(m.east)!0.3!(p.west)$) -- node[left, lbl, text = badRed] {(2) captures} (e);
        \draw[atk] (e) -- node[below right, lbl, text = badRed] {(3) resends the\\same request} (p.south);
        \node[note, right = 0.6cm of p, text width = 3.4cm, draw = okGreen]
        {\textcolor{okGreen}{\textbf{Rejected:}} key \texttt{a81f\ldots} already used; nonce already
            seen; time stamp outside the window};
    \end{tikzpicture}
    \caption{Situation~3: replay attack and its detection.}
    \label{fig:s3}
\end{figure}

\textbf{Тип атаки.} Це атака \emph{повторного відтворення} (replay attack): зловмиснику не потрібно розшифровувати 
чи модифікувати повідомлення достатньо лише повторно надіслати дійсне повідомлення. Навіть запит із дійсним 
підписом приймається, оскільки сам підпис залишається правильним. Наслідок: з клієнта стягується плата двічі або 
більше.

\textbf{Контрзаходи.}
\begin{itemize}[noitemsep]
    \item \textbf{Ключ ідемпотентності (Idempotency key)} -- унікальне випадкове значення (наприклад, UUID v4) для 
        кожної операції. Сервер зберігає результат першого запиту з цим ключем і повертає той самий результат у 
        разі повторення, тому платіж ніколи не виконуватиметься двічі~\cite{stripe_idempotency}.
    \item \textbf{Nonce та мітка часу всередині signed/HMAC-protected part} -- одержувач відхиляє повідомлення, 
        старші за короткий проміжок часу (Stripe використовує приблизно п'ять хвилин для веб-хуків), та запам'ятовує 
        нещодавно отримані nonce значення~\cite{stripe_webhook_signature}.
    \item \textbf{Унікальні ідентифікатори транзакцій та прив'язка до сеансу} -- НБУ вимагає, щоб кожна електронна 
        взаємодія містила унікальний ідентифікатор сеансу, номер операції та мітку часу за UTC~\cite[п.~74]{nbu58}.
    \item \textbf{Одноразові коди автентифікації} -- код 3DS або SCA дійсний лише для однієї операції, тому 
        перехоплений код не може авторизувати інший платіж~\cite[п.~16]{nbu58}.
    \item \textbf{Протокол TLS} з одного боку, запобігає перехопленню запиту в мережі, але не захищає від повторного 
        відтворення запиту з боку скомпрометованого компонента, тому зазначені вище заходи на рівні додатка 
        залишаються необхідними.
\end{itemize}

\subsection{Ситуація 4: Модифікація параметрів платежу}
\label{sec:s4}

Під час здійснення платежу хтось намагається змінити суму або отримувача (Figure~\ref{fig:s4}).

\begin{figure}[H]
    \centering
    \begin{tikzpicture}
        \node[actor] (c) {Customer /\\Merchant};
        \node[evil, right = 2.2cm of c] (e) {Attacker\\(MITM / malware)};
        \node[actor, right = 2.2cm of e] (b) {Bank / PSP};
        \draw[flow] (c) -- node[above, lbl] {600 UAH\\to Shop} (e);
        \draw[atk] (e) -- node[above, lbl, text = badRed] {60\,000 UAH\\to Attacker} (b);
        \node[note, below = 0.6cm of b, text width = 5.6cm, draw = okGreen, xshift = -1.2cm, font = \small]
        {\textcolor{okGreen}{\textbf{Detected:}}
            HMAC/signature over \{amount, payee, order, ts\} does not match;\\
            SCA code is bound to 600~UAH + Shop $\Rightarrow$ invalid;\\
            server recomputes the amount from the order};
    \end{tikzpicture}
    \caption{Situation~4: attempt to modify the amount or payee.}
    \label{fig:s4}
\end{figure}

\textbf{Можливі наслідки.} Клієнт сплачує більше, ніж було домовлено, або гроші надходять на рахунок зловмисника 
замість магазину. Такі атаки можуть здійснюватися за допомогою мережевої атаки man-in-the-middle, шкідливого 
програмного забезпечення в браузері (так званий "man-in-the-browser"{}, що переписує форму), або шляхом фальсифікації 
параметрів на стороні клієнта (наприклад, зміни прихованого поля суми \texttt{amount}).

\textbf{Криптографічні механізми захисту.}
\begin{itemize}[noitemsep]
    \item \textbf{TLS} -- будь-яка зміна даних під час передачі порушує тег автентифікації запису~\cite{rfc9846}.
    \item \textbf{HMAC або цифровий підпис над критичними полями} (сума, валюта, одержувач, ідентифікатор 
        замовлення, мітка часу) -- одержувач перераховує тег; будь-яка зміна виявляється навіть після декількох 
        проміжних вузлів~\cite{rfc2104}.
    \item \textbf{Динамічне зв'язування (Dynamic linking) (НБУ \No.~58)} -- код автентифікації SCA обчислюється 
        на основі суми та одержувача, що відображаються користувачеві; будь-яка зміна робить код недійсним, і 
        операція відхиляється~\cite[п.~22]{nbu58}. Це має вирішальне значення для захисту від шкідливого програмного 
        забезпечення: навіть якщо браузер заражений, користувач бачить справжню суму та одержувача в окремому 
        банківському додатку.
\end{itemize}

\textbf{Механізми на рівні додатків.} Жодної довіри до суми, надісланої клієнтом -- сервер перераховує її на основі 
запиту у своїй базі даних; перевірка одержувача за даними рахунку продавця; реквізити рахунку PSP та реквізити 
виплати можна змінити лише після додаткової автентифікації; ліміти та виявлення аномалій на стороні емітента; у 
системі Open Banking банку заборонено змінювати інструкцію, отриману від ініціатора, а ініціатору заборонено 
змінювати суму~\cite[п.~27, 35]{nbu80}.

\subsection{Ситуація 5: Витік облікових даних користувачів}
\label{sec:s5}

Зловмисник має логін та пароль користувача платіжної системи, але кожна операція вимагає додаткового підтвердження 
(рис~\ref{fig:s5}).

\begin{figure}[H]
    \centering
    \begin{tikzpicture}
        \node[evil] (e) {Attacker\\login + password};
        \node[actor, right = 3cm of e] (s) {Payment\\service};
        \node[actor, right = 3cm of s, draw = okGreen, fill = okGreen!8] (ph) {User's phone\\(possession +\\biometrics)};
        \draw[atk] (e) -- node[above, lbl, text = badRed] {factor 1: OK} (s);
        \draw[flow] (s) -- node[above, lbl] {push / OTP:\\\enquote{600 UAH to X?}} (ph);
        \draw[back, badRed] ([yshift = -4mm]ph.west) -- node[below, lbl, text = badRed] {no confirmation}
        ([yshift = -4mm]s.east);
        \node[note, below = 0.9cm of s, text width = 9.5cm, draw = warnOrange]
        {\textcolor{warnOrange}{\textbf{Residual risks:}} social engineering (\enquote{tell me the code},
            \enquote{press confirm}), SIM swap for SMS OTP, real-time phishing proxy, malware or remote-access
            tool on the phone, push fatigue, abuse of SCA exemptions (small amounts, trusted payees)};
    \end{tikzpicture}
    \caption{Situation~5: stolen password vs.\ multi-factor authentication.}
    \label{fig:s5}
\end{figure}

\textbf{Роль MFA і 3-D Secure.} Пароль є лише \emph{knowledge factor}. Багатофакторна автентифікація (MFA) вимагає 
другого, незалежного фактора -- володіння (\emph{possession}) (зареєстрований телефон, апаратний ключ) або 
притаманності (\emph{inherence}) (відбиток пальця, обличчя)~\cite{nist_80063b4, nbu58}. Зловмисник знає пароль, 
але не має телефону жертви, тому операцію неможливо підтвердити. Для карткових платежів 3-D Secure забезпечує саме 
таку перевірку з боку емітента~\cite{emv_3ds}, а НБУ обмежує кількість спроб введення даних -- п'ятьма невдалими 
спробами~\cite[п.~17]{nbu58}. Крім того, завдяки динамічному зв'язку користувачеві відображаються реальна сума та 
отримувач платежу, тому несподіваний запит на підтвердження є очевидним сигналом тривоги.

\textbf{Залишкові ризики.} Багатофакторна автентифікація (MFA) суттєво знижує, але не усуває (!) ризик:
\begin{itemize}[noitemsep]
    \item \textbf{Соціальна інженерія (social engineering)} -- користувача переконують продиктувати одноразовий 
        пароль (one-time password, OTP) або підтвердити push-повідомлення (дзвінок від "служби безпеки банку"{});
    \item \textbf{Заміна SIM / перехоплення SMS} -- SMS-коди прив'язані до номера телефону, а не до пристрою;
    \item \textbf{Фішинг у реальному часі (зловмисник-посередник)} -- підроблений сайт негайно пересилає 
        одноразовий пароль (OTP); запобігти цьому можуть лише захищені від фішингу засоби автентифікації 
        (FIDO2 / ключі доступу), які прив'язують відповідь до справжнього домену~\cite{nist_80063b4};
    \item \textbf{Зкомпроментований пристрій} -- шкідливе програмне забезпечення або програми віддаленого 
        доступу на телефоні;
    \item \textbf{Винятки} -- платежі на невеликі суми або на користь довірених одержувачів можуть здійснюватися 
        без SCA, згідно до~\cite[п.~37--39]{nbu58}, тому зловмисник може спробувати здійснити багато дрібних 
        платежів або додати себе до списку довірених осіб (що саме по собі вимагає SCA~\cite[п.~38]{nbu58});
    \item \textbf{Витік конфіденційної інформації} -- навіть без здійснення платежу зловмисник може побачити 
        інформацію про обліковий запис, якщо сам факт входу надає доступ до неї.
\end{itemize}
Тому MFA необхідно доповнити моніторингом шахрайства, повідомленнями про кожен вхід та платіж.

\section{Банківський переказ через Open Banking}
\label{sec:openbanking}


\textbf{Open Banking} це структурований та безпечний обмін даними між банком, що веде рахунок, та стороннім постачальником 
послуг через спеціалізовані інтерфейси (APIs), який здійснюється виключно за згодою користувача~\cite[п.~2]{nbu80}. 
В Україні це регулюється Постановою НБУ \No~80 (набула чинності 1 серпня 2025 року) разом із Постановою \No~58~\cite{nbu80, nbu58}. 
Замість введення реквізитів картки клієнт здійснює оплату \emph{безпосередньо з банківського рахунку} шляхом кредитного 
переказу, який ініціює сторонній додаток, а клієнт підтверджує у своєму банку. Українські банки базували специфікацію 
своїх API на європейському стандарті NextGenPSD2 від Berlin Group~\cite{mind_ob_berlin}; на міжнародному рівні найпоширенішим 
профілем безпеки для таких APIs є OpenID Foundation's FAPI (Financial-grade API)~\cite{fapi2}.

\begin{figure}[H]
    \centering
    \begin{tikzpicture}[node distance = 1.2cm]
        \node[actor] (u) {User\\(payer)};
        \node[actor, right = 3.8cm of u] (pisp) {PISP\\(payment initiation\\service provider)};
        \node[actor, right = 3.8cm of pisp] (aspsp) {ASPSP\\(payer's bank)};
        \node[actor, below = 2.3cm of aspsp] (payee) {Payee's bank\\(merchant account)};
        \node[actor, below = 2.3cm of pisp, draw = gray, fill = softGray] (reg) {NBU register +\\qualified trust\\
            service provider};

        \draw[flow] ([yshift = 3mm]u.east) -- node[above, lbl] {(1) \enquote{Pay by bank}\\amount, payee}
        ([yshift = 3mm]pisp.west);
        \draw[flow] ([yshift = 3mm]pisp.east) -- node[above, lbl] {(2) payment request\\mTLS (QWAC)\\+ eSeal (QSealC)} ([yshift = 3mm]aspsp.west);
        \draw[flow, warnOrange] (aspsp.north) to[out = 150, in = 30]
        node[above, lbl, text = warnOrange!70!black, font=\small] {(3) redirect to the bank app: SCA with dynamic linking
            (amount + payee) $\Rightarrow$ consent} (u.north);
        \draw[back] ([yshift = -3mm]aspsp.west) -- node[below, lbl] {(4) auth. code,\\token, status} ([yshift = -3mm]pisp.east);
        \draw[back] ([yshift = -3mm]pisp.west) -- node[below, lbl] {(6) confirmation}
        ([yshift = -3mm]u.east);
        \draw[flow] (aspsp) -- node[right, lbl] {(5) credit transfer\\(SEP / interbank)} (payee);
        \draw[flow, gray, dashed] (aspsp.south west) -- node[above, lbl, sloped] {checks certificate\\and
            authorisation} (reg.east);
    \end{tikzpicture}
    \caption{Спрощена схема платежу, ініційованого через Open Banking}
    \label{fig:openbanking}
\end{figure}

\textbf{Основні учасники}~\cite[п.~3, 6--8]{nbu80}:
\begin{itemize}[noitemsep]
    \item \textbf{Платник (payer)} -- власник рахунку, який надає згоду;
    \item \textbf{ASPSP (Account Servicing Payment Service Provider)} -- постачальник платіжних послуг, що обслуговує 
        рахунок. Фактично це банк, у якому відкрито рахунок; він здійснює автентифікацію користувача та виконує платіж;
    \item \textbf{PISP (Payment Initiation Service Provider)} -- постачальник послуг ініціювання платежів, третя 
        сторона (додаток, каса магазину або інший банк), уповноважена НБУ ініціювати платежі; ПІС (постачальники 
        послуг інформації про рахунки, або англійською -- AISPs (Account Information Service
        Providers)) лише зчитують залишки та історію;
    \item \textbf{Банк одержувача} -- отримує грошовий переказ; НБУ веде реєстр уповноважених постачальників 
        послуг.
    \item \textbf{Технічний оператор} (optional) -- забезпечує технічне з'єднання;
\end{itemize}

\textbf{Дані та криптографічні облікові дані.} Платіжна інструкція (сума, IBAN отримувача, ім'я, призначення 
платежу); кваліфіковані сертифікати PISP для відкритого банкінгу -- сертифікат автентифікації веб-сайту (QWAC), 
що використовується у взаємному TLS, та сертифікат електронної печатки (QSealC), який застосовується для завірення 
платіжного повідомлення~\cite[п.~16, 30--31, 34]{nbu80}; коди авторизації типу OAuth та короткочасні токени доступу, 
прив'язані до сертифіката клієнта~\cite{fapi2, rfc8705}. Власні облікові дані користувача (PIN-код, біометричні 
дані, ключі пристрою) використовуються \emph{виключно} всередині додатка банку.

\textbf{Де зберігається найважливіша інформація.} Сам рахунок, елементи автентифікації користувача та ключі SCA 
залишаються в ASPSP. PISP заборонено зберігати конфіденційні платіжні дані користувача, утримувати кошти користувача 
або змінювати параметри платежу~\cite[п.~35]{nbu80}; він повинен забезпечувати безпеку своїх
власних приватних ключів (для mTLS та електронної печатки(eSeal)), зазвичай у HSM.

\textbf{Як відбувається автентифікація користувача.} Згоду на кожен платіж платник надає банку через власний 
додаток банку з використанням своєї автентифікації, яку повинен здійснювати ASPSP~\cite[п.~47--48]{nbu80}. PISP 
ніколи не бачить банківський пароль користувача — це головна перевага з точки зору безпеки порівняно зі screen scraping 
або введенням реквізитів картки.

\textbf{Як підтверджується операція.} Код автентифікації SCA динамічно пов'язаний із сумою та одержувачем~\cite[п.~22]{nbu58}. 
Після успішного проходження SCA банк видає код авторизації, який PISP обмінює на токен, прив'язаний до відправника; 
потім банк виконує кредитний переказ, не має права змінювати інструкцію та повідомляє PISP про статус~\cite[п.~27--28]{nbu80}.

\textbf{Криптографічні механізми.}
\begin{itemize}[noitemsep]
    \item \textbf{Взаємний (mutual) TLS} із кваліфікованими сертифікатами: банк аутентифікує PISP та перевіряє 
        його у реєстрі НБУ перед кожним доступом~\cite[п.~16]{nbu80};
    \item \textbf{Rваліфікована/розширена електронна печатка} на платіжному повідомленні PISP підтверджує 
        цілісність та невідмовність авторства інструкції~\cite[п.~30]{nbu80};
    \item \textbf{SCA з динамічним зв'язуванням} (одноразові коди або підписи в банківському додатку)~\cite{nbu58};
    \item \textbf{Зміцнення безпеки OAuth~2.0} як у FAPI~2.0: запити на авторизацію, що надсилаються з боку 
        сервера (pushed authorisation requests, PAR), щоб параметри запиту не можна було змінити в браузері~\cite{rfc9126}, 
        PKCE, коди авторизації, дійсні не більше 60 секунд, та токени з обмеженням відправника через mTLS або 
        DPoP, щоб викрадений токен не міг бути використаний іншим клієнтом~\cite{fapi2}.
\end{itemize}

\textbf{Типові загрози.} Фішинг, що імітує перенаправлення на офіційний сайт банку -- користувача переконують підтвердити 
платіж, ініційований шахраєм (шахрайство з авторизованими push-платежами); зловмисний або скомпрометований PISP 
або крадіжка його приватних ключів; перехоплення або витік кодів авторизації та токенів; маніпулювання параметрами 
платежу до SCA; атаки на доступність API (НБУ визначає збій як п'ять запитів, на які не отримано відповіді протягом 
30 секунд~\cite[п.~84]{nbu58}); а також надмірний збір даних третіми сторонами. Оскільки кредитний переказ зазвичай 
є незворотним (на відміну від карткового платежу, тут немає можливості повернення коштів), ретельна перевірка 
користувачем одержувача на екрані SCA має вирішальне значення.

\section{Історичний криптографічний протокол: MicroMint}
\label{sec:micromint}

\textbf{Проблема, що вирішується.} MicroMint був запропонований Рональдом Рівестом та Аді Шаміром у 1996 році 
разом зі схемою PayWord для здійснення \emph{мікроплатежів} -- наприклад, оплата одним центом перегляду однієї 
веб-сторінки~\cite{rivest1996micromint}. У разі карткових платежів комісія за обробку перевищує суму такого 
платежу, а операції з відкритими ключами були дорогими: типова робоча станція того часу могла виконувати 
приблизно два підписи RSA на секунду, але 20\,000 геш-обчислень на секунду~\cite{rivest1996micromint}. Метою 
MicroMint було повністю відмовитися від криптографії з відкритими ключами та звести вартість кожного платежу до 
кількох обчислень геш-функції, погодившись на нижчий, але \textit{достатньо хороший} рівень безпеки -- подібно 
до автоматів з цукерками, що працюють на монети, які не мають дорогих детекторів підробок і все одно функціонують~\cite{rivest1996micromint}.

\textbf{Основна ідея: монета це колізія геш-функції} Учасниками є \emph{broker} (a bank-like mint), \emph{users} 
та \emph{vendors}. Монета -- це набір із $k$ різних значень $x_1,\ldots,x_k$, які дають однаковий геш: 
$h(x_1) = h(x_2) = \ldots = h(x_k)$, де $k$-кратна колізія~\cite{rivest1996micromint}. Створення таких колізій 
схоже на випадкове кидання кульок у контейнери: монета -- це контейнер, до якого потрапило $k$ кульок (Figure~\ref{fig:micromint}). 
Пошук першої колізії є дуже витратним, але після досягнення цього порогу кількість колізій швидко зростає 
($\approx c^k$  монет за $c$ разів більший обсяг роботи). Тому брокер, що володіє потужним спеціалізованим обладнанням, 
виробляє монети дешево, тоді як дрібний фальшивомонетник платить за кожну монету набагато більше, ніж вона коштує~\cite{rivest1996micromint}.

\begin{figure}[H]
    \centering
    \begin{tikzpicture}[scale = 0.9]
        \foreach \i in {0,...,9} {
                \draw[thick] (\i*1.3, 0) -- (\i*1.3, 1.8) -- (\i*1.3 + 1, 1.8) -- (\i*1.3 + 1, 0);
                \node[font = \tiny] at (\i*1.3 + 0.5, -0.25) {bin \i};
            }
        % balls: counts per bin
        \foreach \i/\n in {0/1, 1/4, 2/2, 3/0, 4/3, 5/4, 6/1, 7/2, 8/4, 9/1} {
                \ifnum\n>0
                    \foreach \j in {1,...,\n} {
                            \fill[headBlue!70] (\i*1.3 + 0.5, \j*0.38 - 0.13) circle (0.16);
                        }
                \fi
            }
        \foreach \i in {1, 5, 8} {
                \draw[okGreen, ultra thick, rounded corners] (\i*1.3 - 0.1, -0.05) rectangle (\i*1.3 + 1.1, 1.9);
                \node[font = \scriptsize\bfseries, text = okGreen!70!black] at (\i*1.3 + 0.5, 2.2) {coin};
            }
        \node[note, text width = 12cm, anchor = north, font=\small] at (6.4, -0.6)
        {Ball $x$ goes into bin $y = h(x)$. A bin with $k = 4$ balls gives a coin $(x_1, x_2, x_3, x_4)$.
            The vendor checks a coin with only 4 hash computations; the broker needs $\approx 2^{54}$ hashes
            per month to mint $2^{30}$ coins.};
    \end{tikzpicture}
    \caption{MicroMint: minting coins as $k$-way hash collisions (\enquote{balls into bins}).}
    \label{fig:micromint}
\end{figure}

\textbf{Спрощений алгоритм роботи:}
\begin{enumerate}[noitemsep]
    \item На початку місяця брокер оприлюднює критерій дійсності монет (геш-функцію та необхідні старші біти 
        "правильного"{} значення гешу). Монети дійсні лише протягом цього місяця~\cite{rivest1996micromint}.
    \item Користувачі купують монети у брокера (наприклад, 2~500 монет за 25\$), а брокер фіксує, які монети 
        дісталися якому користувачеві.
    \item Щоб розрахуватися, користувач передає продавцю невикористану монету. Продавець перевіряє її локально: 
        обчислює $k$ гешів і перевіряє, чи вони рівні та "правильні"{} -- зв'язок із брокером не потрібен.
    \item Наприкінці дня продавець повертає зібрані монети брокеру, який оплачує кожну монету не більше ніж один 
        раз і виявляє монети, які були витрачені двічі.
\end{enumerate}
Посередник обирає $k = 4$, інвестує близько \$150\,000 доларів у мікросхеми FPGA та диски й випускає $2^{30}$ 
(близько одного мільярда) монет на місяць, тоді як підробка однієї монети на звичайній робочій станції зайняла 
б приблизно 80 років~\cite{rivest1996micromint}.

\textbf{Криптографічні механізми.} Використовується лише одностороння геш-функція, стійка до колізій (автори 
застосовували гешування на основі DES або MD5/SHA, скорочене до короткого результату). Опціональні розширення 
передбачають використання \emph{hidden predicates} (таємні щоденні умови щодо вартості монет, які розкриваються пізніше) 
та монет, \emph{user-specific} або \emph{vendor-specific}, де структура монети залежить від гешу ідентичності 
користувача або продавця~\cite{rivest1996micromint}.

\textbf{Властивості безпеки.}
\begin{itemize}[noitemsep]
    \item \textbf{Неможливість підробки з економічних міркувань} дрібна підробка коштує дорожче, ніж приносить 
        прибутку; масштабна підробка виявляється, коли з'являються монети з \enquote{unused} контейнерів, і 
        посередник може відкликати серію~\cite{rivest1996micromint};
    \item \textbf{Обмежений термін дії} -- монети втрачають чинність щомісяця, тому фальсифікатор завжди починає 
        з запізненням;
    \item \textbf{double-spending detection} \emph{після} взаємодії з брокером, який ідентифікує постійних шахраїв 
        та виключає їх (незначні шахрайські дії допускаються);
    \item \textbf{Захист від крадіжки.} Монети, прив'язані до конкретного користувача чи продавця, є марними для 
        інших користувачів та вендорів;
    \item \textbf{Відсутність анонімності.} Брокер знає, який користувач купив які монети; MicroMint не є цифровими 
        грошима.
\end{itemize}

\textbf{Comparison with modern payment mechanisms.}

\begin{longtblr}[
        caption = {MicroMint у порівнянні з сучасними механізмами електронних платежів.},
        label = {tab:micromint},
    ]{
        colspec = {Q[l,m,0.2\linewidth] X[l,m] X[l,m]},
        row{1} = {font = \bfseries, bg = headBlue!15},
        rowhead = 1,
        hlines, vlines,
    }
    Аспект & MicroMint (1996) & Сучасні платежі (2026) \\
    Основа безпеки & Економіка колізій хеш-функцій, \enquote{достатньо хороший} рівень безпеки & Надійна криптографія: TLS~1.3, HMAC, підписи, ключі, захищені HSM~\cite{rfc9846, rfc2104} \\
    Перевірка & Офлайн у продавця (4 хеші) & Онлайн-авторизація емітентом для кожного платежу~\cite{stripe_card_process} \\
    Подвійна витрата & Виявляється пізніше; допускаються незначні втрати & Запобігається в режимі онлайн (резерв балансу, ідемпотентні ключі)~\cite{stripe_idempotency} \\
    Прив'язка вартості & Монети, прив'язані до користувача/продавця & Платіжні токени з обмеженням домену (продавець, пристрій, канал)~\cite{emv_tokenisation_qrg} \\
    Обмежений термін дії & Монети втрачають чинність щомісяця & Життєвий цикл токенів, одноразові коди SCA, 60-секундні коди авторизації~\cite{nbu58, fapi2} \\
    Підхід, заснований на оцінці ризиків & \enquote{Не переймайся дрібницями} & Винятки з SCA для платежів з низькою вартістю та низьким рівнем ризику (TRA)~\cite{nbu58} \\
    Автентифікація користувача & Не застосовується до окремих платежів (володіння монетами) & Обов'язкова SCA з динамічним зв'язуванням~\cite{nbu58, emv_3ds} \\
\end{longtblr}

Основна ідея MicroMint -- \emph{зробити кожну перевірку безпеки пропорційною до вартості, що піддається ризику, 
та використовувати дешеві симетричні примітиви замість операцій з відкритим ключем} де це можливо. Це актуально 
і зараз. Сучасні системи використовують HMAC замість підписів для повідомлень з великим обсягом даних (веб-хуки), 
прив'язують токени до домену, як і монети конкретних постачальників, та застосовують винятки на основі ризику 
для дрібних платежів. Різниця лише в тому, що сьогодні дешеві обчислення роблять надійну криптографію доступною 
для кожного платежу, а підключення до Інтернету дозволяє не просто виявляти подвійні витрати, а й запобігати їм. 
Сам MicroMint ніколи не набував широкого поширення: сьогодні питання мікроплатежів зазвичай вирішується шляхом 
агрегації (підписки, передплачені гаманці), а не за допомогою спеціальних монет.

\section{Висновки}
\label{sec:conclusions}

\begin{enumerate}
    \item Онлайн-платіж карткою передбачає участь п'яти сторін -- клієнта, продавця, PSP/еквайра, платіжної 
        мережі та емітента. Аутентифікація ("хто платить?"{}) здійснюється емітентом за допомогою 3-D Secure та 
        SCA, тоді як авторизація ("чи можна виконати цей платіж?"{}) -- це рішення емітента, що ґрунтується на 
        наявності коштів, статусі картки та оцінці ризику шахрайства.
    \item Жоден окремий механізм не забезпечує повного захисту платежу. Протокол TLS захищає дані під час 
        передачі; токенізація та стандарт PCI DSS захищають збережені дані; CVV та 3-D Secure підтверджують 
        володіння карткою та особу її власника; HMAC та цифрові підписи гарантують цілісність та автентичність 
        повідомлень на всіх етапах передачі; ідемпотентні ключі, нонси та часові мітки запобігають повторному 
        відтворенню.
    \item Аналіз п'яти ситуацій порушення безпеки показує, що сам по собі протокол TLS захищає лише від мережевих 
        атак (Ситуація~1); витоки з баз даних запобігаються за допомогою токенізації (Ситуація~2), повторне 
        відтворення -- за допомогою ідемпотентності та одноразових значень (Ситуація~3), модифікація за допомогою 
        автентифікації повідомлень та динамічного зв'язування (Ситуація~4), а крадіжка облікових даних за допомогою 
        багатофакторної автентифікації (Ситуація~5). Залишкові ризики, такі як соціальна інженерія та скомпрометовані 
        пристрої, залишаються і вимагають моніторингу та навчання користувачів.
    \item Українське законодавство (Постанови НБУ \No~58 та \No~4) визначає SCA відповідно до практики ЄС: два 
        незалежні фактори, динамічне прив'язування коду до суми та отримувача, обмеження на кількість невдалих 
    \item Open Banking переносить найважливіші дані та аутентифікацію користувача до його власного банку: 
        сторонній постачальник послуг проходить аутентифікацію за допомогою кваліфікованих сертифікатів, 
        зашифровує свої повідомлення та ніколи не бачить облікових даних користувача. Головною загрозою, що 
        залишається, є соціальна інженерія, оскільки кредитні перекази, як правило, є незворотними.
    \item MicroMint демонструє одну з перших спроб збалансувати вартість безпеки та вартість платежу за 
        допомогою використання виключно геш-функцій. Його основні ідеї це безпека пропорційна вартості, обмежений 
        термін дії та прив'язка до користувача/продавця -- знову з'являються у сучасній токенізації, короткострокових 
        облікових даних та винятках із SCA.
\end{enumerate}